%%%%%%%%%%%%%%%%%%%%%%%%%%%%%%%%%%%%%%%%%%%%%%%%%%%%%%%%%%%%%%%%%%%%%%%%
%% Art. 27, VI — "se necessário, utilizar exemplos e/ou quadros        %%
%% comparativos, relacionando-os com o estado da técnica".              %%
%%                                                                     %%
%% Resultados de teste comparativo costumam ser decisivos para          %%
%% demonstrar o efeito técnico, e por isso pesam no exame. Se você tem  %%
%% esses dados, este é o lugar deles.                                   %%
%%                                                                     %%
%% TABELAS aqui, FIGURAS nunca:                                         %%
%%   art. 20 — tabelas, fórmulas ou estruturas químicas e expressões    %%
%%     matemáticas podem ser inseridas no texto, em preto e            %%
%%     identificadas de forma sequencial. Use o ambiente 'tabelainpi'   %%
%%     (numera e rotula sozinho) e o ambiente 'equation' para           %%
%%     expressões matemáticas;                                          %%
%%   art. 19 — representações gráficas (figuras, fotografias,          %%
%%     fluxogramas, gráficos) NÃO são aceitas no relatório descritivo.  %%
%%     Elas vão exclusivamente no documento de desenhos.                %%
%%                                                                     %%
%% Toda tabela precisa ser comentada no texto — uma tabela solta, sem   %%
%% parágrafo que a explique, é matéria que o examinador não sabe como   %%
%% ler.                                                                 %%
%%%%%%%%%%%%%%%%%%%%%%%%%%%%%%%%%%%%%%%%%%%%%%%%%%%%%%%%%%%%%%%%%%%%%%%%

\secaoinpi{Exemplos de concretizações \danatureza}

\ifinvencao
\pnum Em uma concretização preferida, o corpo tubular (2) é obtido por injeção
de poliamida reforçada com fibra de vidro, a membrana elástica (3) é de
elastômero de etileno-propileno-dieno e a mola de retorno (5) é de aço inoxidável
austenítico, com pré-carga de 12 N.
\else
\pnum Em uma primeira variação construtiva, o corpo tubular (2) é obtido por
injeção de poliamida reforçada com fibra de vidro, a membrana elástica (3) é de
elastômero de etileno-propileno-dieno e a mola de retorno (5) é de aço inoxidável
austenítico, com pré-carga de 12 N.
\fi

\pnum O atuador eletromecânico (6) é constituído por um solenoide de bobina
única, alimentado em 12 V, e o sensor de pressão (7) é do tipo piezorresistivo,
com faixa de medição de 0 kPa a 600 kPa.

\pnum A energia consumida por ciclo de acionamento é obtida pela integral da
potência aplicada ao atuador eletromecânico (6) ao longo da duração do pulso,
conforme a expressão:

\begin{equation}
    W = \int_{0}^{t_{p}} v(t)\, i(t)\, \mathrm{d}t
\end{equation}

\pnum Foram realizados ensaios comparativos entre uma concretização do
dispositivo de acionamento (1) e dois conjuntos representativos do estado da
técnica, submetidos ao mesmo regime de 40 ciclos diários de irrigação, com
pressão de linha de 200 kPa. Os resultados estão reunidos na Tabela 1.

\begin{tabelainpi}{Energia consumida por ciclo de acionamento e pressão mínima
de fechamento}
    \begin{tabular}{lrr}
        \toprule
        Conjunto ensaiado & Energia por ciclo (J) & Pressão mínima (kPa) \\
        \midrule
        Acionamento direto (BR 10 2015 000000 0)      & 216{,}0 & 0   \\
        Acionamento hidráulico (US 2018/0123456 A1)   & 4{,}8   & 120 \\
        Dispositivo de acionamento (1)                & 1{,}9   & 0   \\
        \bottomrule
    \end{tabular}
\end{tabelainpi}

\pnum Os valores reunidos na Tabela 1 mostram que o dispositivo de acionamento
(1) consumiu 1,9 J por ciclo, contra 216,0 J do conjunto de acionamento direto
e 4,8 J do conjunto de acionamento hidráulico assistido. O conjunto de
acionamento hidráulico assistido deixou de completar o fechamento com pressão
de linha inferior a 120 kPa, enquanto o dispositivo de acionamento (1) completou
o fechamento em todos os ensaios, inclusive com a linha despressurizada.
